% 2009 年新课标卷理科第 19 题：四棱锥 S-ABCD，P 为侧棱 SD 上一点。
% 来源 PDF：普通高考/2009/2009新课标理(海南,宁夏).pdf，
% SHA256 f42b575218180e261a89cf8733109d4960b61560a97b5f266731a51c422a4045
% 原页：第 2 页（A3 横向）中栏中部，印刷页 542。
% 重绘依据：底面 ABCD 为正方形，S 在底面中心正上方；P 在 SD 上。
% 虚线为底面被遮挡的边 AC、BD 及锥高 SO。
\begin{tikzpicture}[scale=0.92, semithick, >=Stealth,
  every node/.style={inner sep=1.5pt, font=\small}]
  \coordinate (A) at (-2.4,-1.1);
  \coordinate (B) at (2.0,-1.1);
  \coordinate (C) at (2.8,0.7);
  \coordinate (D) at (-1.6,0.7);
  \coordinate (S) at (0.2,3.9);
  \coordinate (O) at (0.2,-0.2);
  \coordinate (P) at (-0.79,2.13);
  % 底面
  \draw (A) -- (B) -- (C) -- (D) -- cycle;
  \draw[dashed] (A) -- (C);
  \draw[dashed] (B) -- (D);
  % 侧棱
  \draw (S) -- (A);
  \draw (S) -- (B);
  \draw (S) -- (C);
  \draw (S) -- (D);
  % 锥高与 P
  \draw[dashed] (S) -- (O);
  \fill (P) circle (2.1pt) node[above left] {$P$};
  \draw[dashed] (A) -- (P);
  \draw[dashed] (C) -- (P);
  % 顶点标注
  \node[below left]  at (A) {$A$};
  \node[below right] at (B) {$B$};
  \node[right]       at (C) {$C$};
  \node[left]        at (D) {$D$};
  \node[above]       at (S) {$S$};
  \node[below]       at (O) {$O$};
\end{tikzpicture}
